\documentclass[10pt,a4paper]{amsart}
\usepackage{amsmath,amssymb,amsthm,graphicx}
% Retained minimal preamble for the Stacks Project chapter more-morphisms.tex (More on Morphisms).
% Verbatim excerpt of the chapter preamble: Xy-pic setup (native preamble lines 14-18) and the native
% theorem declarations with the mathematical macros (native preamble lines 154-217). Copied mechanically;
% nothing added, removed or redefined. The wrapper class and the amsmath/amssymb/amsthm/graphicx base
% packages are supplied by the assembler, so the class/externaldocument part of the native preamble is omitted.
% For commutative diagrams we use Xy-pic
\usepackage[all]{xy}

% We use 2cell for 2-commutative diagrams.
\xyoption{2cell}
% Theorem environments.
%
\theoremstyle{plain}
\newtheorem{theorem}[subsection]{Theorem}
\newtheorem{proposition}[subsection]{Proposition}
\newtheorem{lemma}[subsection]{Lemma}

\theoremstyle{definition}
\newtheorem{definition}[subsection]{Definition}
\newtheorem{example}[subsection]{Example}
\newtheorem{exercise}[subsection]{Exercise}
\newtheorem{situation}[subsection]{Situation}

\theoremstyle{remark}
\newtheorem{remark}[subsection]{Remark}
\newtheorem{remarks}[subsection]{Remarks}

\numberwithin{equation}{subsection}

% Macros
%
\def\lim{\mathop{\mathrm{lim}}\nolimits}
\def\colim{\mathop{\mathrm{colim}}\nolimits}
\def\Spec{\mathop{\mathrm{Spec}}}
\def\Hom{\mathop{\mathrm{Hom}}\nolimits}
\def\Ext{\mathop{\mathrm{Ext}}\nolimits}
\def\SheafHom{\mathop{\mathcal{H}\!\mathit{om}}\nolimits}
\def\SheafExt{\mathop{\mathcal{E}\!\mathit{xt}}\nolimits}
\def\Sch{\mathit{Sch}}
\def\Mor{\mathop{\mathrm{Mor}}\nolimits}
\def\Ob{\mathop{\mathrm{Ob}}\nolimits}
\def\Sh{\mathop{\mathit{Sh}}\nolimits}
\def\NL{\mathop{N\!L}\nolimits}
\def\CH{\mathop{\mathrm{CH}}\nolimits}
\def\proetale{{pro\text{-}\acute{e}tale}}
\def\etale{{\acute{e}tale}}
\def\QCoh{\mathit{QCoh}}
\def\Ker{\mathop{\mathrm{Ker}}}
\def\Im{\mathop{\mathrm{Im}}}
\def\Coker{\mathop{\mathrm{Coker}}}
\def\Coim{\mathop{\mathrm{Coim}}}

% Boxtimes
%
\DeclareMathSymbol{\boxtimes}{\mathbin}{AMSa}{"02}

%
% Macros for moduli stacks/spaces
%
\def\QCohstack{\mathcal{QC}\!\mathit{oh}}
\def\Cohstack{\mathcal{C}\!\mathit{oh}}
\def\Spacesstack{\mathcal{S}\!\mathit{paces}}
\def\Quotfunctor{\mathrm{Quot}}
\def\Hilbfunctor{\mathrm{Hilb}}
\def\Curvesstack{\mathcal{C}\!\mathit{urves}}
\def\Polarizedstack{\mathcal{P}\!\mathit{olarized}}
\def\Complexesstack{\mathcal{C}\!\mathit{omplexes}}
% \Pic is the operator that assigns to X its picard group, usage \Pic(X)
% \Picardstack_{X/B} denotes the Picard stack of X over B
% \Picardfunctor_{X/B} denotes the Picard functor of X over B
\def\Pic{\mathop{\mathrm{Pic}}\nolimits}
\def\Picardstack{\mathcal{P}\!\mathit{ic}}
\def\Picardfunctor{\mathrm{Pic}}
\def\Deformationcategory{\mathcal{D}\!\mathit{ef}}
\newenvironment{reference}{}{}
\newenvironment{slogan}{}{}
\newenvironment{history}{}{}
\usepackage{cleveref}
\crefname{originallemma}{original lemma}{original lemmas}
\Crefname{originallemma}{Original lemma}{Original lemmas}
\makeatletter
\newlabel{descent-lemma-locally-finite-nr-irred-local-fppf}{{0BAL}{source}{}{sourcepointer-1}{}}
\newlabel{descent-lemma-locally-finite-nr-irred-local-fppf@cref}{{[originallemma][0][]0BAL}{[1][0][]source}{}{}{}}
\makeatother
\crefname{originallemma}{original lemma}{original lemmas}
\Crefname{originallemma}{Original lemma}{Original lemmas}
\makeatletter
\newlabel{descent-lemma-normal-local-smooth}{{034F}{source}{}{sourcepointer-2}{}}
\newlabel{descent-lemma-normal-local-smooth@cref}{{[originallemma][0][]034F}{[1][0][]source}{}{}{}}
\makeatother
\crefname{originallemma}{original lemma}{original lemmas}
\Crefname{originallemma}{Original lemma}{Original lemmas}
\makeatletter
\newlabel{descent-lemma-reduced-local-smooth}{{034E}{source}{}{sourcepointer-3}{}}
\newlabel{descent-lemma-reduced-local-smooth@cref}{{[originallemma][0][]034E}{[1][0][]source}{}{}{}}
\makeatother
\crefname{originallemma}{original lemma}{original lemmas}
\Crefname{originallemma}{Original lemma}{Original lemmas}
\makeatletter
\newlabel{morphisms-lemma-normalization-normal}{{035Q}{source}{}{sourcepointer-4}{}}
\newlabel{morphisms-lemma-normalization-normal@cref}{{[originallemma][0][]035Q}{[1][0][]source}{}{}{}}
\makeatother
\title{Normalization commutes with a smooth morphism under a finiteness condition on components}
\author{The Stacks Project Authors}
\date{stacks-project:more-morphisms.tex:local-snapshot-20261011; GNU FDL 1.2 or later}
\begin{document}
\maketitle
\noindent\textit{Editorial note.} Statement span is lines 5409-5415 of sources/stacks-more-morphisms/pinned-original-source.tex; proof span is lines 5417-5436, the complete author proof (18 body lines).\par\medskip
\noindent\textit{Editorial note.} Navigation stubs: 4 cited-but-not-included labels are supplied as non-load-bearing pointers with their pinned tags where the snapshot tag index has them.\par\medskip
\noindent\textit{Editorial note.} extra\_packages is empty on purpose: the retained preamble already loads xy with options and a second declaration clashes in the assembler wrapper.\par\medskip
\noindent\textit{Editorial note.} Text is retained verbatim, including the author's own typos and placeholders.\par\medskip
\noindent\textit{Allowed standard prerequisites (editorial).}\begin{itemize}
\item Explicit prerequisite boundary. Normalizations, smooth morphisms, and irreducible components of quasi-compact opens; the comparison is retained.
\end{itemize}\medskip
\ifdefined\hypertarget\hypertarget{sourcepointer-1}{}\fi
\noindent\textit{Source reference.} Original lemma 0BAL: sources/stacks-more-morphisms/pinned-original-source.tex, source lines 5436--5436. Cited by the retained statement or proof but not part of the collected spans; supplied only as a navigation pointer. The argument does not depend on its content.\par\medskip
\ifdefined\hypertarget\hypertarget{sourcepointer-2}{}\fi
\noindent\textit{Source reference.} Original lemma 034F: sources/stacks-more-morphisms/pinned-original-source.tex, source lines 5436--5436. Cited by the retained statement or proof but not part of the collected spans; supplied only as a navigation pointer. The argument does not depend on its content.\par\medskip
\ifdefined\hypertarget\hypertarget{sourcepointer-3}{}\fi
\noindent\textit{Source reference.} Original lemma 034E: sources/stacks-more-morphisms/pinned-original-source.tex, source lines 5436--5436. Cited by the retained statement or proof but not part of the collected spans; supplied only as a navigation pointer. The argument does not depend on its content.\par\medskip
\ifdefined\hypertarget\hypertarget{sourcepointer-4}{}\fi
\noindent\textit{Source reference.} Original lemma 035Q: sources/stacks-more-morphisms/pinned-original-source.tex, source lines 5436--5436. Cited by the retained statement or proof but not part of the collected spans; supplied only as a navigation pointer. The argument does not depend on its content.\par\medskip
\begin{lemma}[Normalization and smooth morphisms]
\label{lemma-normalization-and-smooth}
Let $X \to Y$ be a smooth morphism of schemes. Assume every quasi-compact
open of $Y$ has finitely many irreducible components. Then the same
is true for $X$ and there is a unique isomorphism $X^\nu = X \times_Y Y^\nu$
over $X$ where $X^\nu$, $Y^\nu$ are the normalizations of $X$, $Y$.
\end{lemma}

\begin{proof}
By Descent, Lemma \ref{descent-lemma-locally-finite-nr-irred-local-fppf}
every quasi-compact open of $X$ has finitely many irreducible components.
Note that $X_{red} = X \times_Y Y_{red}$ as a scheme smooth over a reduced
scheme is reduced, see
Descent, Lemma \ref{descent-lemma-reduced-local-smooth}.
Hence we may assume that $X$ and $Y$ are reduced (as the normalization
of a scheme is equal to the normalization of its reduction by
definition). Next, note that $X' = X \times_Y Y^\nu$ is a normal scheme
by Descent, Lemma \ref{descent-lemma-normal-local-smooth}.
The morphism $X' \to Y^\nu$ is smooth (hence flat) thus the generic
points of irreducible components of $X'$ lie over generic points of
irreducible components of $Y^\nu$. Since $Y^\nu \to Y$ is birational
we conclude that $X' \to X$ is birational too (because $X' \to Y^\nu$
induces an isomorphism on fibres over generic points of $Y$).
We conclude that there exists a factorization
$X^\nu \to X' \to X$, see
Morphisms, Lemma \ref{morphisms-lemma-normalization-normal}
which is an isomorphism as $X'$ is normal and integral over $X$.
\end{proof}

\end{document}
